\subsection{映射的锥}

设 X 与 Y 为拓扑空间，f 为从 X 到 Y 中的连续映射。设 \(\operatorname{Cyl}(f)\) 为映射 f 的柱。用 \(\alpha_{f}\) 表示从 \(X \times I\) 到 \(\operatorname{Cyl}(f)\) 中的典范映射，用 \(f_{0}\) 表示映射 \(x \mapsto \alpha_{f}(x,0)\)（从 X 到 \(\operatorname{Cyl}(f)\) 中）；这些映射分别诱导出 \(X \times I\) 与 X 到它们在 \(\operatorname{Cyl}(f)\) 中的像上的同胚。

定义 11. --- 映射 f 的锥记作 \(\operatorname{Côn}(f)\)，它是由 \(\operatorname{Cyl}(f)\) 缩塌 \(f_{0}(\mathrm{X})\) 所得的拓扑空间。

用 \(\beta_{f}^{\prime}\colon\mathrm{Y}\to\mathrm{C}\hat{\mathrm{o}}\mathrm{n}(f)\) 表示典范映射 \(\beta_{f}\colon\mathrm{Y}\to\mathrm{C}\mathrm{y}\mathrm{l}(f)\) 与从 \(\mathrm{C}\mathrm{y}\mathrm{l}(f)\) 到 \(\mathrm{C}\hat{\mathrm{o}}\mathrm{n}(f)\) 上的典范满射的复合。映射 \(\beta_{f}^{\prime}\) 是连续的，并定义了从 Y 到 \(\mathrm{C}\hat{\mathrm{o}}\mathrm{n}(f)\) 的一个闭子集上的同胚，该闭子集称为锥底，我们据此把它与 Y 等同。

用 \(\alpha_{f}^{\prime}\colon\mathrm{X}\times\mathbf{I}\to\mathrm{C}\hat{\mathrm{o}}\mathrm{n}(f)\) 表示映射 \(\alpha_{f}\) 与从 \(\mathrm{Cyl}(f)\) 到 \(\mathrm{C}\hat{\mathrm{o}}\mathrm{n}(f)\) 上的典范满射的复合。映射 \(\alpha_{f}^{\prime}\) 是一个同伦，其初始映射是常值映射，末项映射是映射 \(\beta_{f}^{\prime}\circ f\)。

若 X 为空，则映射 \(\beta_{f}^{\prime}\) 是从 Y 到 \(\operatorname{Côn}(f)\) 上的同胚。

假设 X 非空。用 s 表示空间 \(\mathrm{Cyl}(f)/f_{0}(\mathrm{X})\) 的基点；它称为锥 \(\mathrm{Côn}(f)\) 的锥顶点。由于 \(f_{0}(\mathrm{X})\) 在 \(\mathrm{Cyl}(f)\) 中闭，典范映射 \(\pi\colon\mathrm{Cyl}(f)\to\mathrm{Côn}(f)\) 诱导出 \(\mathrm{Cyl}(f)\setminus f_{0}(\mathrm{X})\) 到 \(\mathrm{Côn}(f)\setminus\{s\}\) 上的同胚（III, p. 252）。锥 \(\mathrm{Côn}(f)\) 的锥顶点不属于其锥底；Y 到 \(\mathrm{Côn}(f)\setminus\{s\}\) 中的典范单射是同伦等价（III, p. 239, 命题 6）。

设 \(\sigma_f\colon \mathrm{Cyl}(f)\times \mathbf{I}\to \mathrm{Cyl}(f)\) 为 \(f\) 的柱到其柱底上的典范收缩。对 \(c\in \mathrm{Cyl}(f)\setminus f_0(\mathrm{X})\) 与 \(t\in \mathbf{I}\)，有 \(\sigma_f(c,t)\notin f_0(\mathrm{X})\)。因此，存在唯一的映射 \(\sigma_f'\colon (\mathrm{Côn}(f)\setminus\{s\})\times \mathbf{I}\to \mathrm{Côn}(f)\setminus\{s\}\)，满足 \(\sigma_f'(\pi(c),t)=\pi(\sigma_f(c,t))\)（对 \(c\in \mathrm{Cyl}(f)\setminus f_0(\mathrm{X})\) 与 \(t\in \mathbf{I}\)）。它是连续的，并且是 \(\mathrm{Côn}(f)\setminus\{s\}\) 到 Y 上的一个强收缩。它称为 \(\mathrm{Côn}(f)\setminus\{s\}\) 到其锥底上的典范收缩。它的末项映射是 \(\mathrm{Côn}(f)\setminus\{s\}\) 到 Y 上的一个收缩映射，称为去掉锥顶点的锥到其锥底上的典范收缩映射。

命题 11（锥的泛性质）

设 Z 为拓扑空间，\(\beta: \mathrm{Y} \to \mathrm{Z}\) 为连续映射，\(\alpha: \mathrm{X} \times \mathbf{I} \to \mathrm{Z}\) 为一同伦，其初始映射是常值映射，末项映射等于 \(\beta \circ f\)。存在唯一的连续映射 \(\varphi\)，它从 \(\operatorname{Côn}(f)\) 到 Z 中，满足 \(\alpha = \varphi \circ \alpha_f'\) 且 \(\beta = \varphi \circ \beta_f'\)。

由柱的泛性质（III, p. 238, 命题 4），存在唯一的连续映射 \(h \colon \operatorname{Cyl}(f) \to \mathbf{Z}\)，满足 \(\alpha = h \circ \alpha_f\) 且 \(\beta = h \circ \beta_f\)。由于 \(\alpha\) 的初始映射是常值映射，\(h\) 在子空间 \(\alpha_f(\mathbf{X} \times \{0\})\) 上的限制是常值的。因此存在唯一的连续映射 \(\varphi \colon \operatorname{Côn}(f) \to \mathbf{Z}\)，满足 \(h = \varphi \circ \pi\)，其中 \(\pi\) 表示从 \(\operatorname{Cyl}(f)\) 到 \(\operatorname{Côn}(f)\) 上的典范满射。映射 \(\varphi\) 满足 \(\alpha = \varphi \circ \alpha_f'\) 与 \(\beta = \varphi \circ \beta_f'\)，并且是唯一具有这些性质的映射，因为 \(\alpha_f'\) 与 \(\beta_f'\) 的像覆盖 \(\operatorname{Côn}(f)\)。

例 1. --- 设 X 为拓扑空间。我们把映射 IdX 的锥称为空间 X 的锥，记作 C(X)；它是由 X × I 缩塌 X × \{0\} 所得的空间。设 \(\pi\) 为从 X × I 到 C(X) 中的典范映射；X × {[}0, 1{[} 在 \(\pi\) 下的像 C'(X) 称为空间 X 的开锥。

假设 X 非空。那么锥 C(X) 非空，其基点仍称为锥的顶点。

映射 \(((x,t),u)\mapsto(x,t(1-u))\)（从 \((\mathrm{X}\times\mathbf{I})\times\mathbf{I}\) 到 \(X\times I\) 中）是连接 \(X\times I\) 的恒同映射与映射 \((x,t)\mapsto(x,0)\) 的同伦。由此，通过取商，我们得到映射 \(\sigma\colon\mathrm{C}(\mathrm{X})\times\mathbf{I}\to\mathrm{C}(\mathrm{X})\)，满足 \(\sigma(\pi(x,t),u)=\pi(x,t(1-u))\)（对 \(x\in X\)，\(t,u\in I\)）。由 I, p. 19, 命题 10，映射 \(\sigma\) 是连续的。于是映射 \(\sigma\) 是 s 处的带基点同伦，连接 \(\mathrm{C}(\mathrm{X})\) 的恒同映射与以 \(\{s\}\) 为像的常值映射。因此，锥 \(\mathrm{C}(\mathrm{X})\) 在其顶点 s 处可缩。由于 \(\sigma(\mathrm{C}'(\mathrm{X})\times\mathbf{I})\) 包含于 \(\mathrm{C}'(\mathrm{X})\)，映射 \(\sigma\) 通过取子空间定义了 s 处的带基点同伦，连接 \(\mathrm{C}'(\mathrm{X})\) 的恒同映射与以 \(\{s\}\) 为像的常值映射；因此开锥 \(\mathrm{C}'(\mathrm{X})\) 在 s 处可缩。

注记。--- 1) 设 X 为拓扑空间，A 为 X 的子空间；用 i 表示 A 到 X 中的典范单射。典范收缩映射 \(r_{i}\)（它把柱 \(\mathrm{Cyl}(i)\) 收缩到其柱底上）把子空间 \(\alpha_{i}(\mathrm{A} \times \{0\})\) 映入 A 中。设 \(\rho: \mathrm{C}\hat{\mathrm{o}}\mathrm{n}(i) \to \mathrm{X}/\mathrm{A}\) 为由此通过取商诱导的连续映射。假设偶 (X, A) 具有同伦扩张性质。由 III, p. 243, 定理 1，映射 \(r_{i}\) 定义了偶 \((\mathrm{Cyl}(i), \alpha_{i}(\mathrm{A} \times \{0\}))\) 到偶 (X, A) 上的同伦等价。由 III, p. 253 的注 4，映射 \(\rho\) 便是偶 \((\mathrm{C}\hat{\mathrm{o}}\mathrm{n}(i), s)\) 到偶 \((\mathrm{X}, s_{\mathrm{X}/\mathrm{A}})\) 上的同伦等价。特别地，它是一个同伦等价。

\begin{enumerate}
\def\labelenumi{\arabic{enumi})}
\setcounter{enumi}{1}
\tightlist
\item
  设 X 与 Y 为拓扑空间，\(f\colon X \to Y\) 为连续映射。典范映射 \(\alpha_{f}^{\prime}\colon X \times I \to \operatorname{Côn}(f)\) 在 \(X \times \{0\}\) 上是常值的，因此定义了连续映射 \(\gamma_{f}\colon C(X) \to \operatorname{Côn}(f)\)。限制到 \(C'(X)\) 上时，映射 \(\gamma_{f}\) 是单射且严格的，并通过取子空间诱导出开锥 \(C'(X)\) 到 Y 在 \(\operatorname{Côn}(f)\) 中的补集上的同胚。
\end{enumerate}

把锥 \(C(X)\) 的锥底与空间 X 等同，并用 \(u\) 表示从 \(Y \cup_f C(X)\) 到 \(\operatorname{Côn}(f)\) 中的、由映射 \(\beta_f':Y\to\operatorname{Côn}(f)\) 与 \(\gamma_f:C(X)\to\operatorname{Côn}(f)\) 诱导的连续映射。反之，设 \(v:\operatorname{Côn}(f)\to Y\cup_f C(X)\) 为唯一的连续映射，满足 \(v\circ\alpha_f'\) 是从 \(X\times I\) 到 \(C(X)\) 上的典范映射，\(v\circ\beta_f'\) 是从 \(Y\) 到 \(Y\cup_f C(X)\) 的典范映射。映射 \(u\) 与 \(v\) 是互逆的同胚。

